%%%PAGE 118 rot=0 legibility=good summary=曲线运动(轨迹与合力方向)、电场中带电粒子轨迹分析(电场线/等势面/能量与电势比较)、点电荷等势面、异种电荷轨迹
%%%BLOCK id=118-01 ch=04 sec=曲线运动·轨迹与合力方向 kind=knowledge alt=09 cont=none y=0.05-0.18
% 页面右上角露出邻页(化学)内容“混合气体 ρ=1.429g/L … 22.4 … 6.72L”，不属于本页，已忽略
\textbf{曲线运动}

速度合力夹轨迹\quad 凹侧方向偏向力

切线$v$\quad 凹侧力\quad 锐加钝减
%%%END
%%%BLOCK id=118-02 ch=09 sec=带电粒子在电场中的轨迹分析 kind=method alt=04 cont=none y=0.14-0.36
电场线\quad 切$+$凹

等势面\quad 垂$+$凹\quad 处处垂$E$，疏密不同

%FIG p118_a 0.10 0.195 0.36 0.35
\notefig[0.3]{figures/p118_a.png}
% 图中标注：A、B 两点，F，多条带箭头的曲线轨迹

\[
\left\{\begin{aligned}
E_{kB}&>E_{kA}\\
E_{pB}&<E_{pA}\ \text{力方向}
\end{aligned}\right.
\qquad
\begin{aligned}
&\text{电性\ 正同负反}\\
&W_{\text{电}}\ \text{正同负反}
\end{aligned}
\]

\[
\left\{\begin{aligned}
a_A&>a_B\\
\varphi_A&>\varphi_B
\end{aligned}\right.
\ \text{疏密、顺逆电场线}
\qquad
\text{沿受力方向\ 动增势减}
\]
%%%END
%%%BLOCK id=118-03 ch=09 sec=点电荷电场·等势面 kind=knowledge alt=none cont=none y=0.35-0.43
%FIG p118_b 0.125 0.35 0.24 0.43
\notefig[0.25]{figures/p118_b.png}

圆形辐状电场等势面越往外越宽，非等间隔（\unclear{$\pm Q$}）
%%%END
%%%BLOCK id=118-04 ch=09 sec=带电粒子在电场中的轨迹分析 kind=method alt=none cont=none y=0.44-0.57
%FIG p118_c 0.13 0.445 0.25 0.568
\notefig[0.25]{figures/p118_c.png}
% 图中标注：弧线轨迹，弧上带 $\oplus$ 的点，左侧一个小电荷符号，受力 F 箭头向左（指向凹侧）

异种电荷

凹侧包住点电荷、离的近时动能大
%%%END
%%%BLOCK id=118-05 ch=09 sec=带电粒子在电场中的轨迹分析 kind=method alt=none cont=none y=0.56-0.65
%FIG p118_d 0.14 0.562 0.32 0.648
\notefig[0.25]{figures/p118_d.png}
% 图中标注：左侧孤立的 $\oplus$，两条相对弯曲的弧线（弧上各有一个 $\oplus$），左弧上有一个 X 记号

粒子轨迹从X到X可以不确定
%%%END
%%%PAGEEND 118
